\documentclass{article}
\usepackage{amsmath, amssymb, amsthm}
\usepackage{hyperref}
\usepackage{geometry}
\geometry{margin=1in}

\title{Erd\H{o}s Problem \#51}
\date{Last updated: 2025-08-31}
\author{Source: erdosproblems.com}

\begin{document}
\maketitle

\noindent\textbf{Status:} OPEN \quad \textbf{No prize}

\noindent\textbf{Tags:} number theory

\medskip

\noindent\textbf{Problem Statement:}

Is there an infinite set $A\subset \mathbb{N}$ such that for every $a\in A$ there is an integer $n$ such that $\phi(n)=a$, and yet if $n_a$ is the smallest such integer then $n_a/a\to \infty$ as $a\to\infty$?




Carmichael has asked whether there is an integer $t$ for which $\phi(n)=t$ has exactly one solution. Erd\H{o}s has proved that if such a $t$ exists then there must be infinitely many such $t$. See also [694] . This is discussed in problems B36 and B39 of Guy's collection \cite{Gu04}.




References

[Gu04] Guy, Richard K., Unsolved problems in number theory . (2004), xviii+437.

\noindent\textbf{Related OEIS sequences:} A002202, A014197


\bigskip
\noindent\small{Source: \url{https://www.erdosproblems.com/51}}
\end{document}
